\chapter{分布式训练与大规模工程：并行策略、训练框架、成本与排障}

\begin{introduction}
\item {\heiti 本章核心问题}：当企业知识助手背后的训练、继续预训练与大规模实验继续扩张时，团队应该如何选择并行策略、训练框架和排障方法，才能让系统既跑得动，也管得住？
\item {\heiti 主案例推进}：把企业知识助手从“单机可实验”推进到“多卡可训练、多任务可治理、故障可恢复”的平台期工程。
\item {\heiti 读完后能做什么}：面对单卡放不下、GPU 利用率上不去、通信开销过大、检查点过慢、实验难复现等问题时，知道该先用哪类策略、再查哪一层瓶颈。
\end{introduction}

\section{学习目标}
\begin{itemize}
  \item 理解数据并行、张量并行、流水线并行、参数分片与 3D 并行各自主要解决什么瓶颈。
  \item 掌握 DDP、FSDP、DeepSpeed、Accelerate 与 Megatron 一类技术栈的适配边界。
  \item 建立分布式训练里的容错、检查点、成本核算与实验复现意识。
  \item 能根据显存瓶颈、通信瓶颈、I/O 瓶颈和调度复杂度做出更稳的工程取舍。
\end{itemize}

\section{问题背景}
如果说前面的章节都还可以在“单机工程”的心智模型里理解，那么走到这里以后，事情已经变了。团队不再只关心模型会不会答、LoRA 能不能训，而要同时关心：多张 GPU 是否真的带来更高吞吐；一次失败的训练能否从中断点恢复；日志、超参数和数据版本能否对齐；一次实验到底花掉了多少 GPU 小时。

这正是很多人第一次真正感受到“大模型工程”和“普通深度学习实验”之间的差别。单机实验里，最常见的问题是脚本报错、显存溢出或者 loss 不降；分布式工程里，更常见的问题却是通信等待、阶段不均衡、首卡显存过高、检查点写盘太慢、恢复训练后状态不一致、同一个实验下周无法复现。

所以，本章要回答的并不是“并行技术有哪些”，而是一个更务实的问题：\textbf{当模型规模变大后，系统复杂度如何被重新分配到计算、显存、通信、存储和治理之中。} 只有把这件事讲清楚，前面的数据工程、训练、评测和推理优化，才会在平台期继续成立。

\section{核心概念}
\subsection{分布式训练首先是资源分账，而不是多卡魔法}
单机训练时，很多资源问题被隐藏在一台机器内部；分布式训练则迫使团队把问题拆开记账。粗略地说，训练资源至少可以拆成五类：
\begin{enumerate}
  \item \textbf{参数存储}：模型权重本身要放在哪里。
  \item \textbf{训练状态}：梯度、优化器状态、激活值和临时缓冲区如何分担。
  \item \textbf{通信成本}：哪些张量需要同步、谁和谁同步、同步频率多高。
  \item \textbf{失败恢复}：训练中断后如何尽量少损失进度。
  \item \textbf{治理成本}：实验记录、版本对齐、资源统计和故障归因是否可持续。
\end{enumerate}

分布式训练真正难的地方，不是概念多，而是任何一种并行策略都在重新分配这五类成本。你得到某一项收益，往往也会付出另一项代价。

\subsection{先看瓶颈属于哪一类，再谈并行策略}
旧稿里把很多方案按技术名称展开，但从教材视角看，团队更需要一张“瓶颈到策略”的映射表：
\begin{table}[H]
\centering
\small
\begin{tabularx}{\textwidth}{>{\raggedright\arraybackslash}p{3cm}>{\raggedright\arraybackslash}p{3.2cm}>{\raggedright\arraybackslash}X}
\toprule
\textbf{当前瓶颈} & \textbf{优先策略} & \textbf{原因} \\
\midrule
单卡装不下，但单层算子还不算大 & FSDP / ZeRO / 参数分片 & 先解决常驻显存问题，复杂度相对可控 \\
单层矩阵太大，单卡连一层都放不下 & 张量并行 TP & 这是典型的层内切分场景 \\
模型很深，层间分配更自然 & 流水线并行 PP & 通过分阶段前后向换取可训练性 \\
多卡算得动，但同步太慢 & 先查 DDP/FSDP 配置、桶大小、网络与 batch & 这通常是通信和调度问题，不是模型问题 \\
多种瓶颈同时出现，且已有成熟集群 & 3D 并行组合 & 只在规模和基础设施都足够时再上复杂组合 \\
\bottomrule
\end{tabularx}
\caption{并行策略不是从名字选，而是从瓶颈选}
\label{tab:distributed-bottleneck-routing}
\end{table}

\subsection{通信基础：点对点与集体通信}
在所有并行策略背后，最底层的共同问题都是通信。教材里最值得保留的通信划分只有两个：
\begin{itemize}
  \item \textbf{点对点通信}：更常见于流水线并行，相邻阶段之间传递激活和梯度。
  \item \textbf{集体通信}：更常见于数据并行和张量并行，例如 Broadcast、Reduce、All-Reduce、Reduce-Scatter、All-Gather。
\end{itemize}

这两个词的价值在于，它们直接决定你该怀疑什么。如果是点对点通信，重点看阶段切分是否均衡；如果是集体通信，重点看同步频率、张量大小和网络拓扑是否在吞收益。

\begin{table}[H]
\centering
\small
\begin{tabularx}{\textwidth}{>{\raggedright\arraybackslash}p{2.8cm}>{\raggedright\arraybackslash}p{3.4cm}>{\raggedright\arraybackslash}X}
\toprule
\textbf{通信类型} & \textbf{典型出现位置} & \textbf{最该关注的工程问题} \\
\midrule
点对点通信 & 流水线并行 PP & 相邻阶段负载是否均衡、是否出现等待和气泡 \\
集体通信 & DDP、FSDP、TP & All-Reduce / All-Gather 是否频繁到吞掉扩卡收益 \\
\bottomrule
\end{tabularx}
\caption{不同并行策略背后对应不同通信心智模型}
\label{tab:distributed-communication}
\end{table}

\subsection{Ring-AllReduce 的直觉：为什么扩卡收益会先撞到同步墙}
分布式训练里最容易被误解的一点，是“多卡同步”并不是一个抽象黑盒。对数据并行来说，局部前反向算完以后，各卡必须把梯度汇成同一个全局结果；Ring-AllReduce 之所以经典，就是因为它在大规模 GPU 集群里比“所有卡都把数据扔给一台主卡”更可扩展。

工程上不必背完算法细节，只要抓住这条直觉：\textbf{每增加一张卡，你得到的是更多本地计算资源，同时也引入了更多同步对象。} 于是扩卡收益什么时候开始变差，往往不只取决于模型算得多快，也取决于梯度多大、同步多频繁、网络多强，以及 batch 是否大到足以摊薄同步成本。

把它说得更落地一点，Ring-AllReduce 往往会让团队在下面三件事上最早吃亏：
\begin{itemize}
  \item \textbf{batch 太小}：每步本地计算太少，通信时间很难被摊薄。
  \item \textbf{梯度碎片太多}：同步调用频繁，通信开销被切得很碎。
  \item \textbf{跨机网络太弱}：本地卡虽然多了，但全局最慢链路开始主导节奏。
\end{itemize}

这也是为什么扩卡失败时，第一反应不该只是“模型实现有问题”，而应先问：\textbf{我现在是在算得太少、同步得太多，还是网络本身跟不上。}

\subsection{数据并行：最常见的起点，但不是万能起点}
数据并行的基本思路最容易理解：把同一个模型复制到多张卡上，每张卡吃不同 batch，再同步梯度。它常常是团队进入多卡训练的第一站，也是很多中等规模任务最实用的基线。

但教材里一定要明确一点：\textbf{数据并行解决的是吞吐扩展问题，不解决“单卡装不下完整模型”的根矛盾。} 如果每张卡都必须保留完整权重、完整梯度和完整优化器状态，那么显存瓶颈很快就会出现。

\subsection{为什么 DDP 通常比 DataParallel 更值得讲}
旧稿里专门讲了 \texttt{nn.DataParallel} 和 \texttt{DistributedDataParallel}。教材里不需要保留太多实现细节，但必须保留结论：\textbf{现代多卡训练默认应优先 DDP，而不是 DataParallel。}

原因很直接：
\begin{itemize}
  \item DDP 采用多进程方式，避免单进程多线程带来的额外瓶颈。
  \item DDP 的梯度同步更贴近真实的分布式训练路径。
  \item 它既适用于单机多卡，也更自然扩展到多机多卡。
\end{itemize}

如果一个团队还在用 DataParallel 作为长期多卡方案，通常说明它还停留在“多卡试跑”而不是“分布式工程”阶段。

\subsection{FSDP 与 ZeRO：很多时候先分片，比先上复杂并行更稳}
当问题主要是“模型和训练状态太大，单卡放不下”，比起直接切张量或切流水线，很多团队先从 FSDP 或 ZeRO 一类参数分片方案入手更稳。它们的共同思路是：不要让每张卡都完整持有参数、梯度和优化器状态，而是把这些东西分散存储。

\begin{table}[H]
\centering
\small
\begin{tabularx}{\textwidth}{>{\raggedright\arraybackslash}p{2.5cm}>{\raggedright\arraybackslash}p{3.2cm}>{\raggedright\arraybackslash}p{2.6cm}>{\raggedright\arraybackslash}X}
\toprule
\textbf{方案} & \textbf{主要解决什么} & \textbf{主要收益} & \textbf{主要代价} \\
\midrule
DDP & 多卡吞吐扩展 & 实现简单，生态成熟 & 每卡仍要持完整模型与大部分训练状态 \\
FSDP & 参数和状态分片 & 降单卡显存，贴近原生 PyTorch & 配置理解门槛更高 \\
ZeRO-1/2/3 & 逐步分片优化器、梯度、参数 & 显存收益逐步增强 & 通信、配置和恢复逻辑更复杂 \\
\bottomrule
\end{tabularx}
\caption{先做状态分片，往往比直接上重并行更符合中等规模工程现实}
\label{tab:ddp-fsdp-zero}
\end{table}

从决策顺序看，如果团队只是“单卡训练顶到了显存天花板”，通常先怀疑是不是该上分片，而不是立刻上 TP 或 PP。

\subsection{ZeRO 阶段怎么选：先问你最缺哪一类显存}
旧稿里把 ZeRO-1、2、3 分得很细。教材版里最重要的不是记名字，而是记住：\textbf{三个 stage 对应的是三类不同的“冗余状态”被逐步分出去。} 这也意味着，stage 越往后，显存收益越大，通信、恢复与配置复杂度也越高。

\begin{table}[H]
\centering
\small
\begin{tabularx}{\textwidth}{>{\raggedright\arraybackslash}p{2.4cm}>{\raggedright\arraybackslash}p{3.4cm}>{\raggedright\arraybackslash}p{2.8cm}>{\raggedright\arraybackslash}X}
\toprule
\textbf{ZeRO 阶段} & \textbf{主要分掉什么} & \textbf{更适合什么处境} & \textbf{最该提醒什么} \\
\midrule
Stage 1 & 优化器状态 & 模型还能放下，但优化器状态开始顶显存 & 改造相对轻，但解决不了参数本体过大问题 \\
Stage 2 & 优化器状态 + 梯度 & 反向阶段显存明显吃紧，且希望继续走数据并行路径 & 通信压力会上升，恢复链也要跟着更完整 \\
Stage 3 & 优化器状态 + 梯度 + 参数 & 模型本体已经逼近或超过单卡能力上限 & 显存收益最大，但通信、配置、保存与恢复也最复杂 \\
\bottomrule
\end{tabularx}
\caption{ZeRO 的选择不只是“能不能省显存”，而是“愿不愿意接住对应复杂度”}
\label{tab:zero-stage-selection}
\end{table}

对大多数工程团队来说，一个很稳的顺序是：\textbf{先把 DDP 路径跑明白，再试 Stage 1/2，只有当参数本体真的装不下时，才认真评估 Stage 3。} 这样做的重点不是保守，而是让每一次复杂度升级都对应一个清楚的瓶颈。

\subsection{AMP 与混合精度：先把每一步训练成本降下来}
旧稿里还反复强调过 AMP（Automatic Mixed Precision）和混合精度训练，这一点在新版里也不能丢。很多团队一提“大规模训练”就先想到并行策略，但在真正跨到更复杂的分布式方案之前，\textbf{先把单步训练的显存和算力效率做好}往往更划算。混合精度做的事情并不神秘：让一部分计算与存储用更低精度完成，从而减少显存占用并提升吞吐。

对工程团队来说，最重要的不是背公式，而是记住三条判断：
\begin{itemize}
  \item \textbf{它是训练基线，不只是优化技巧}：如果硬件支持 bf16 或 fp16，通常应该先把混合精度纳入默认实验配置，再讨论是否需要更重的并行策略。
  \item \textbf{bf16 和 fp16 不是一回事}：bf16 动态范围更友好，很多现代训练任务里更稳；fp16 往往更依赖 loss scaling，一旦处理不好，就更容易出现梯度下溢、loss 抖动或训练不收敛。
  \item \textbf{混合精度解决的是“每步成本”，不是“所有容量问题”}：它能延后显存瓶颈到来，但如果模型和优化器状态本身已经远超单卡容量，最终还是要回到 FSDP、ZeRO 或更复杂的并行方案。
\end{itemize}

如果把决策顺序拉直，可以这样理解：\textbf{先看是否能靠混合精度、梯度累积、合理 batch 和检查点策略把单卡或轻量多卡跑顺；只有这些基础手段仍然不够，才升级到更重的分片与并行。} 这也是为什么很多旧稿里的“显存优化技巧”虽然看起来零碎，其实都在服务同一个目标：尽量把复杂度留到真正需要的时候。

\subsection{loss scaling 与溢出监控：fp16 能跑，不等于 fp16 跑稳}
旧稿里专门讲过动态损失缩放，这里值得保留，因为它正好补上了混合精度最容易被误解的一层。很多团队知道“开 fp16 会更省显存”，却不知道 \textbf{fp16 训练真正的脆点往往不是前向能不能跑，而是反向梯度会不会因为数值范围不足而下溢或溢出。}

这也是为什么 fp16 路线通常要配合 loss scaling。直觉上，它是在反向传播前把损失放大，让小梯度先别太早被冲成 0；更新前再把梯度缩回原尺度。动态损失缩放进一步做的是：一旦检测到溢出，就回退缩放倍率；若一段时间都稳定，再尝试把倍率往上调。工程上不一定要手写这套逻辑，但必须知道它在解决什么。

对教材读者来说，更值得记住的是诊断顺序：如果混合精度一开，loss 开始偶发爆炸、梯度出现 \texttt{inf/nan}、某些 step 被频繁跳过，先别急着怀疑数据脏了，也别急着把并行策略全换掉。\textbf{先看当前是 bf16 还是 fp16，再看 loss scaling、梯度裁剪和优化器状态是否在一起工作。} 这一步在平台期尤其关键，因为一旦团队把混合精度写进默认配置，它就不再是一个“可选加速器”，而是训练可靠性的一部分。

\subsection{Offload 是兜底，不是免费的第二层显存}
旧稿里还讲过 ZeRO Offload、CPU Offload 甚至 NVMe Offload。这些内容如果完全删掉，读者很容易误以为“显存不够就把东西丢到 CPU 去”是一条几乎没有代价的路。实际并不是这样。{\heiti offload 的本质，是用更慢的带宽和更高的时延去换一部分显存空间。}

这件事之所以重要，是因为训练状态并不只有模型权重，还包括梯度、优化器状态、激活相关中间结果。把其中一部分移出 GPU 确实能让大模型“先跑起来”，但也会带来新的瓶颈：
\begin{itemize}
  \item \textbf{CPU Offload}：比纯 GPU 常驻更省显存，但 PCIe 带宽和主机内存访问时延会拖慢训练节奏。
  \item \textbf{NVMe Offload}：进一步换空间，适合“实在放不下、先求能跑”的场景，但速度影响往往更明显，对存储带宽和稳定性也更敏感。
  \item \textbf{恢复与排障复杂度上升}：状态不再全在 GPU 内部，checkpoint、恢复、吞吐波动与资源监控都会更复杂。
\end{itemize}

所以更稳的工程判断不是“能不能 offload”，而是“当前目标到底是\textbf{先验证可行性}，还是\textbf{追求稳定高吞吐}”。如果只是想验证更大模型或更长序列是否能训通，offload 是很现实的兜底；但如果团队已经进入稳定迭代阶段，却还长期依赖重度 offload 才能训练，那么往往说明硬件配置、并行策略或任务边界需要重新审视。

\subsection{张量并行：解决的是单层算子太大}
张量并行把一个层里的矩阵运算切到多张卡上，适合模型宽度很大、单层算子本身就装不下或算不过来的场景。它的工程价值在于解决\textbf{层内切分}问题，而不是一般意义上的“多卡就更高级”。

它也对硬件提出更苛刻的要求。因为同一层的计算被拆到了多卡之间，通信会更频繁、更紧密，所以更依赖同机高速互联。若网络条件一般，TP 很容易从“理论上高效”变成“实际中卡在同步上”。

\subsection{TP 内部常见切法：行并行与列并行只影响通信位置}
旧稿里还专门拆过 row parallel 和 column parallel。教材里不需要把每个矩阵维度都画出来，但要给读者留一个足够用的判断：\textbf{张量并行不是笼统地“把层切开”，而是要决定把哪一侧切开，以及通信发生在前向还是后向的哪个位置。}

以全连接层为例，常见的两种思路可以粗略理解为：
\begin{itemize}
  \item \textbf{列并行}：先把输出维度切开，各卡各算一部分输出，后续再做拼接或归并。
  \item \textbf{行并行}：先把输入相关的乘法分散到各卡，最后把局部结果求和或规约回来。
\end{itemize}

这两种切法的核心差别，不在名字，而在\textbf{通信发生在哪里、与哪些算子耦合得最紧。} 也正因为如此，Megatron 一类框架的价值之一，就是帮团队把这些通信位置设计成可复用模板，而不是每次自己重写一遍层内并行逻辑。

\subsection{流水线并行：用调度复杂性换可训练性}
流水线并行把不同层分配给不同设备，让前向和反向像装配线一样流动。它最适合模型很深、层间切分自然的场景。它的重要工程代价是\textbf{气泡}：并不是每张卡每个时刻都在满载计算。

旧稿里关于 GPipe、PipeDream 和 1F1B 的细节很多，教材里只保留一个够用的判断：
\begin{itemize}
  \item \textbf{GPipe 类思路}：更容易理解，但阶段等待更明显。
  \item \textbf{1F1B / PipeDream 类思路}：通过更积极的调度减少空转，但实现、状态管理和调试都更复杂。
\end{itemize}

气泡的存在意味着：流水线并行不是“白送的扩展”，而是用更复杂的调度换来更大模型的可训练性。若团队还没有稳定的监控与排障能力，PP 不能轻易当成默认解。

\subsection{micro-batch 数量会直接决定流水线气泡}
很多团队知道流水线并行有气泡，却没有把它和 micro-batch 数量直接连起来。一个够用的工程直觉是：\textbf{micro-batch 太少，流水线填不满；micro-batch 太多，调度和重算开销又会上来。} 因此，PP 的调参从来不只是“切几层”，还包括“每次把多少个小批次送进这条流水线”。

教材里可以保留一个近似判断：
\[
\text{bubble ratio} \approx \frac{\text{PP stages} - 1}{\text{micro-batches}}.
\]
它不是严格推导公式，但足够说明趋势：阶段数固定时，micro-batch 越少，空转占比越高；micro-batch 增加后，气泡会下降，但不会无限免费，因为 kernel launch、调度、激活保存与重算压力也会跟着涨。

这也是为什么 1F1B 一类调度虽然常被说成“更先进”，但它解决的只是{\heiti 减少等待与缓解内存压力}，不是免除了 micro-batch 选择。真正上线时，团队仍要观察 step time、气泡占比、激活显存和吞吐曲线，找到“气泡开始明显下降，但额外调度成本还没反噬”的那一段，而不是只凭经验固定一个数字。

\subsection{3D 并行不是毕业礼，而是最后的组合拳}
当模型规模继续扩大，团队可能需要把数据并行、张量并行、流水线并行组合起来。此时常见的表达是 3D 并行，其总设备数可以写成：
\[
\text{world size} = \text{DP} \times \text{TP} \times \text{PP}.
\]

这条公式本身不难，难的是它背后的组织成本。因为一旦三类并行同时出现，通信组、日志、故障恢复、性能分析和 checkpoint 恢复都会一起变复杂。所以 3D 并行不该被当作“更高级”的默认路线，而应被看作“规模逼到这里了，且团队基础设施已经准备好了”的结果。

\begin{figure}[H]
\centering
\begin{tikzpicture}[
  font=\small,
  box/.style={llm box, minimum width=2.9cm, minimum height=1cm},
  arr/.style={llm arr}
]
\node[box] (dp) {数据并行 DP\\复制模型\\切 batch};
\node[box, right=0.9cm of dp] (tp) {张量并行 TP\\切矩阵\\解层内瓶颈};
\node[box, right=0.9cm of tp] (pp) {流水线并行 PP\\切层\\解深度瓶颈};
\node[box, below=1.15cm of tp] (zero) {FSDP / ZeRO\\分片参数、梯度、优化器状态};
\draw[arr] (dp) -- (tp);
\draw[arr] (tp) -- (pp);
\draw[arr] (tp) -- (zero);
\end{tikzpicture}
\caption{大规模训练中常见的并行与分片策略关系}
\label{fig:distributed-strategies}
\end{figure}

\subsection{拓扑先于卡数：NVLink、PCIe 与跨机网络会决定扩卡收益}
旧稿在通信章节里其实一直在提醒一件很容易被忽视的事：\textbf{决定扩卡收益上限的，往往不是你有几张卡，而是这些卡之间怎么连。} 对训练工程来说，单机 NVLink、高速 PCIe、跨机以太网或 InfiniBand，并不是背景配置，而是并行策略成立与否的前提条件。

把这件事讲直白一点：
\begin{itemize}
  \item \textbf{数据并行更怕跨机同步慢}：All-Reduce 一旦被网络拖住，卡再多也只是一起等。
  \item \textbf{张量并行更怕同机互联差}：同层算子要频繁交换张量，没有高速互联时，很容易算得比传得慢。
  \item \textbf{流水线并行更怕阶段不均衡}：某一段卡得最慢，整条流水线就跟着最慢段走。
\end{itemize}

所以“先上哪种并行”从来不只由模型决定，也由机房拓扑决定。若当前只有跨机网络较弱的多节点资源，很多时候先把单机内的 FSDP、ZeRO 或混合精度打满，比仓促上 TP/PP 更现实；反过来，若单机内有很强的 GPU 互联，而单层算子已经装不下，那么 TP 的性价比会明显提高。旧稿里那些看似偏底层的通信基础内容，本质上是在帮读者建立这层硬件直觉。

\subsection{框架选型不是问“谁最强”，而是问“谁最适合当前阶段”}
从工程视角看，训练框架选型至少要回答四个问题：
\begin{enumerate}
  \item 现有代码和模型生态是否兼容。
  \item 团队是否能调试、复现和维护。
  \item 当前瓶颈需要哪类并行或分片能力。
  \item 框架带来的收益，是否值得它带来的复杂度。
\end{enumerate}

\begin{table}[H]
\centering
\small
\begin{tabularx}{\textwidth}{>{\raggedright\arraybackslash}p{2.6cm}>{\raggedright\arraybackslash}p{3.1cm}>{\raggedright\arraybackslash}p{2.8cm}>{\raggedright\arraybackslash}X}
\toprule
\textbf{框架 / 技术栈} & \textbf{更适合什么阶段} & \textbf{主要优势} & \textbf{主要提醒} \\
\midrule
Accelerate & 快速原型、轻量多卡训练 & 上手快，屏蔽很多分布式细节 & 不等于极限规模最优方案 \\
DDP / FSDP & 通用多卡训练与原生 PyTorch 生态 & 路径清晰、易与现有代码整合 & 需要团队自己理解更多细节 \\
DeepSpeed & 大模型训练、ZeRO 与复杂优化需求 & 分片与优化能力成熟、可扩展性强 & 配置项多，调试成本高 \\
Megatron 一类技术栈 & TP / PP / 3D 并行需求明确时 & 面向大规模并行优化深入 & 工程门槛和迁移成本更高 \\
\bottomrule
\end{tabularx}
\caption{框架选型的核心不是先进性，而是阶段适配度}
\label{tab:framework-selection}
\end{table}

\subsection{框架迁移成本经常比第一次跑通更贵}
教材里如果只保留“哪个框架适合哪个阶段”，还不够，因为平台期团队真正会被拖慢的，往往不是第一次把脚本跑起来，而是\textbf{后续迁移、维护和统一治理的成本。} 例如，一个团队今天用 Accelerate 快速起了原型，明天切到原生 FSDP，后天又为了某个 ZeRO 特性改投 DeepSpeed。单次看都合理，连起来却很容易形成一条谁也不完全掌握的多栈链路。

更稳的做法，通常不是永远不迁移，而是把迁移当成一次正式工程动作。至少要先回答三件事：第一，当前迁移是在解决明确瓶颈，还是只是在追更“更强”的名词；第二，迁移后谁来维护启动脚本、配置模板、日志接口和 checkpoint 兼容；第三，旧实验资产还能不能被新栈接住。如果这三件事说不清，框架升级很可能会带来一段时间的性能收益，却把更长时间的排障和复现成本留给团队。

\begin{table}[H]
\centering
\small
\begin{tabularx}{\textwidth}{>{\raggedright\arraybackslash}p{3cm}>{\raggedright\arraybackslash}p{3cm}>{\raggedright\arraybackslash}X}
\toprule
\textbf{迁移动机} & \textbf{什么时候算合理} & \textbf{最该提前确认什么} \\
\midrule
从单机栈迁到 DDP / FSDP & 单卡或轻量多卡已经稳定，但显存或吞吐明显碰顶 & 启动脚本、日志和 checkpoint 口径是否能统一迁移 \\
从原生栈迁到 DeepSpeed & ZeRO、offload 或更复杂优化已经是明确需求 & 团队是否愿意接住更重的配置与恢复复杂度 \\
从通用栈迁到 TP / PP / Megatron 路线 & 单层算子或模型深度已逼到原路径上限 & 当前拓扑、模型结构和运维能力是否真的匹配 \\
\bottomrule
\end{tabularx}
\caption{框架迁移的关键不是“能不能切”，而是“切完之后谁来长期维护”}
\label{tab:framework-migration-cost}
\end{table}

\subsection{全局 batch 是一笔乘法账：micro-batch、梯度累积与 DP 必须一起看}
旧稿在实践建议里多次提到 batch size 缩放规则，这里需要在教材版里讲透。很多分布式训练问题看起来像“学习率不对”或“多卡不稳”，实际根源是团队没有把全局 batch 当成一笔完整账来算。对大多数数据并行训练，可以先记住一个最常用的近似关系：
\[
\text{global batch} \approx \text{micro-batch per device} \times \text{gradient accumulation} \times \text{DP replicas}.
\]

这条式子看起来简单，但它决定了很多关键动作。首先，\textbf{梯度累积是在用时间换显存}。单卡吃不下更大的 micro-batch 时，可以通过多走几步再同步，换取更大的全局 batch。其次，\textbf{学习率、warmup 步数和日志口径都要跟着全局 batch 一起重算}。如果只把卡数翻倍，却不重算有效 batch，训练曲线和吞吐比较就会一起失真。最后，\textbf{TP 和 PP 通常不直接增加样本并发数}，它们更多是在解决“单个样本训不下”或“单层算不过来”的问题，所以不要把所有并行都误当成同一种扩 batch 手段。

对团队协作而言，最稳的做法通常是把四个数字固定写进实验记录：每卡 micro-batch、梯度累积步数、数据并行副本数、推导出的全局 batch。只要这笔账清楚，很多“同样代码为什么两次结果不一样”的问题会当场缩小一半。

\subsection{梯度累积和梯度检查点常一起出现，但它们不是一回事}
在大模型训练里，这两个词几乎总是成对出现，于是初学者很容易把它们混成一个“省显存开关”。实际上，它们解决的是两件不同的事：
\begin{itemize}
  \item \textbf{梯度累积}：通过多次前反向再同步，换来更大的有效全局 batch，本质上是在用更多 step 时间换更大 batch。
  \item \textbf{梯度检查点}：不保存部分中间激活，而在反向时重算，本质上是在用更多重计算换更低激活显存。
\end{itemize}

两者经常一起开，是因为团队同时面临两种压力：一方面单卡装不下想要的 micro-batch，另一方面全局 batch 又不能太小。但要注意，这种组合虽然常见，却往往意味着\textbf{单步时间显著上升}。如果只看到“终于不 OOM 了”，却不重新记录每步耗时、每秒 token 和通信占比，就会误以为训练只是“慢一点”，实际上可能已经把有效吞吐打到很低。

更稳的做法是把它们分别记账：梯度累积决定的是同步频率和有效 batch；梯度检查点决定的是激活显存和重算成本。只有把这两条线分开，团队才能回答清楚“现在是显存不够，还是吞吐已经不值得了”。

\subsection{先用小模型和小数据把链路跑通，再放大}
旧稿里有一条很朴素但极重要的经验：\textbf{先用更小的模型、更短的数据、更便宜的配置把训练链路跑通，再把规模放大。} 这不是“保守”，而是分布式排障最省钱的顺序。因为很多问题和模型是否足够大根本无关，而是出在启动脚本、数据管线、混合精度、checkpoint、日志出口或通信配置上。

如果一开始就把所有复杂度叠到最大，团队很容易面对一个无法分层定位的故障现场：到底是数据错、环境错、NCCL 错、优化器状态错，还是模型本身太大？相反，只要先在小模型上验证单机、再验证单机多卡、再验证多机，排障路径就会清楚得多，平台期治理也更容易形成固定动作。

\subsection{启动纪律比“能跑一次”更重要：torchrun、DeepSpeed 与环境一致性}
旧稿里有不少看起来像“命令行细节”的内容，例如 \texttt{torchrun} 怎么起、多机参数怎么传、环境变量怎么对齐。这些内容在教材化时不能简单删掉，因为它们直接决定分布式训练到底是\textbf{可重复启动的工程系统}，还是“运气好时能跑起来一次的脚本”。

对多数 PyTorch 生态团队来说，{\heiti 启动脚本应该被当成配置入口，而不是随手拼接的命令行。} 无论是 \texttt{torchrun} 还是 DeepSpeed 风格的 launcher，最重要的是把节点数、rank、主节点地址、端口、可见 GPU、日志目录和 checkpoint 目录都收束到统一配置里，而不是散落在不同 shell 命令、手工导出环境变量和临时注释里。

\begin{table}[H]
\centering
\small
\begin{tabularx}{\textwidth}{>{\raggedright\arraybackslash}p{3cm}>{\raggedright\arraybackslash}p{3.5cm}>{\raggedright\arraybackslash}X}
\toprule
\textbf{启动层面} & \textbf{最低要求} & \textbf{为什么重要} \\
\midrule
进程编排 & 明确 world size、node rank、local rank 与设备映射 & 避免“程序在跑，但各卡角色不一致”这类隐蔽错误 \\
主节点信息 & 固定 MASTER\_ADDR、MASTER\_PORT 或对应 launcher 配置 & 减少多机初始化失败、端口冲突与随机卡死 \\
设备可见性 & 统一 CUDA 设备可见性、容器设备映射与节点资源声明 & 避免 rank 与 GPU 绑定关系混乱 \\
通信环境 & 明确 NCCL 相关配置、网络接口与超时策略 & 多机不稳定时，启动环境往往比训练脚本更先出问题 \\
输出目录 & 日志、checkpoint、配置快照按实验唯一标识落盘 & 没有统一出口，就没有可靠恢复与复现 \\
\bottomrule
\end{tabularx}
\caption{分布式训练启动不只是“起进程”，而是把运行环境显式写清楚}
\label{tab:launch-discipline}
\end{table}

进一步说，平台期团队最好形成一条很朴素的纪律：\textbf{训练命令只接受少量高层参数，底层环境由统一脚本或配置生成。} 这样当你从单机多卡扩到多机、从 DDP 扩到 FSDP 或 DeepSpeed 时，变化的是配置层，而不是每次都重新手写一遍命令。旧稿里那些关于环境变量、端口、节点拓扑和启动参数的提醒，本质上讲的就是这件事。

\subsection{NCCL 与多机启动最小检查表}
多机训练一旦起不来，很多团队会直接怀疑脚本或框架；但旧稿里的经验更实用：\textbf{先把通信环境当成独立系统排查。} 对多数 NVIDIA GPU 集群来说，这一步通常绕不开 NCCL 与网络配置。

\begin{table}[H]
\centering
\small
\begin{tabularx}{\textwidth}{>{\raggedright\arraybackslash}p{3cm}>{\raggedright\arraybackslash}p{3.4cm}>{\raggedright\arraybackslash}X}
\toprule
\textbf{检查项} & \textbf{主要想确认什么} & \textbf{如果出问题常会表现成什么} \\
\midrule
MASTER\_ADDR / MASTER\_PORT & 所有节点是否真的连到同一主节点与端口 & 初始化随机卡死、部分 rank 永远等不到其他节点 \\
NCCL 网卡接口配置 & 通信是否走了正确网卡 & 明明网络可用，但吞吐异常低或偶发超时 \\
NCCL\_DEBUG 与日志级别 & 是否能拿到足够故障线索 & 出问题时只有“卡住了”，没有可定位证据 \\
IB / P2P 开关与拓扑识别 & 当前环境是否真的支持相应高速链路 & 理论上该快的配置，实际却退化成慢路径 \\
超时与重试策略 & 慢节点或短暂抖动是否会直接把任务打死 & 集群偶发波动就导致整场实验重启 \\
\bottomrule
\end{tabularx}
\caption{排查多机启动问题时，先把通信环境看清，再怀疑训练脚本}
\label{tab:nccl-checklist}
\end{table}

这张表的价值不在于让读者死记环境变量，而在于建立顺序感：\textbf{先看能不能连、走哪张网卡、有没有日志，再谈框架参数。} 没有这层顺序，很多多机问题会在错误层面上反复兜圈。

\subsection{数据加载是隐藏的第四瓶颈：DistributedSampler、预处理与落盘速度}
旧稿虽然没有把这一点单独拎成标题，但很多实践提醒都在指向同一件事：\textbf{训练慢，不一定慢在模型和通信，也可能慢在数据。} 当团队开始多卡、多机训练后，数据加载常常会从一个“不起眼的本地脚本”变成真正的系统瓶颈。

这类瓶颈最常见的形态有三种。第一，\textbf{样本没有按 rank 正确切分}，不同进程重复读到同一批数据，既浪费吞吐，也污染统计口径。第二，\textbf{预处理过重}，例如在线做复杂切分、OCR、清洗或 JSON 解析，导致 GPU 在等 CPU 和磁盘。第三，\textbf{checkpoint 与日志落盘抢带宽}，训练本身不慢，但每隔一段时间就被 I/O 拖住。

因此平台期最好把数据层也当成分布式系统治理：
\begin{itemize}
  \item 明确使用 \texttt{DistributedSampler} 或等价机制，保证每个 rank 的样本分片可解释、可复现；
  \item 尽量把重预处理前移到离线阶段，训练时只保留轻量拼装；
  \item 把 dataloader 吞吐、CPU 利用率、磁盘与网络 I/O 一并纳入训练观测，而不是只看 GPU 面板。
\end{itemize}

这部分一旦做好，很多“GPU 看起来没满”“多卡扩上去却没更快”的问题会立刻变得更好诊断。因为你终于能分清：到底是模型算得慢，还是数据没跟上。

\subsection{环境栈治理：底层库升级要比 Python 包更谨慎}
旧稿里还有一条非常现实的提醒：分布式训练环境里，底层库升级的风险往往比普通 Python 包大得多。驱动、CUDA、NCCL、GLIBC、容器基础镜像和网络栈一旦改动，问题通常不会体现在“导入失败”，而会体现在多机初始化卡死、通信异常、性能突然掉线或某些节点行为不一致。

因此，平台期团队最好把环境栈当成受控资产治理，而不是谁机器坏了就顺手升一下版本。更稳的做法通常是：固定基础镜像和驱动组合；先在小规模回归环境验证；记录每次升级对应的训练脚本版本与压测结果；高风险底层库不要和大规模训练试验同时改。对训练工程来说，\textbf{环境稳定性本身就是吞吐和复现能力的一部分。}

\subsection{检查点、弹性恢复与自动重启不是“运维附属项”}
旧稿里的训练经验提醒非常值得保留：大模型训练常常不是几个小时，而是几天、几周甚至更久。只要训练足够长，中断就不是“会不会发生”，而是“什么时候发生”。因此平台期团队至少要建立三件基础设施：
\begin{itemize}
  \item \textbf{检查点策略}：多长时间保存一次，保什么，保到哪里。
  \item \textbf{恢复策略}：恢复后是否能连同优化器状态、学习率调度和数据进度一起回来。
  \item \textbf{自动重启与弹性容错}：节点失败后能否自动拉起，尽量减少人工盯守。
\end{itemize}

这些工作看似不提升模型精度，却直接决定实验是否可持续。因为训练越贵，中断一次的损失就越大。

\subsection{检查点不只存模型：恢复粒度决定你是在“接着跑”还是“重开局”}
很多初学者以为“能把权重存下来”就算有 checkpoint 了，但平台期训练真正需要的是\textbf{恢复粒度}。如果只存模型参数，下次往往只能重新组织优化器状态、学习率调度和数据进度；这更像“从旧权重重开一局”，而不是“从上次中断处继续跑”。

\begin{table}[H]
\centering
\small
\begin{tabularx}{\textwidth}{>{\raggedright\arraybackslash}p{2.7cm}>{\raggedright\arraybackslash}p{3.2cm}>{\raggedright\arraybackslash}p{2.8cm}>{\raggedright\arraybackslash}X}
\toprule
\textbf{保存粒度} & \textbf{通常包含什么} & \textbf{更适合什么场景} & \textbf{主要限制} \\
\midrule
仅模型权重 & 参数本体 & 推理发布、阶段性导出、做下游微调起点 & 不能完整恢复训练过程 \\
完整训练状态 & 参数、优化器、调度器、随机种子、数据进度 & 长时间训练、继续预训练、昂贵实验 & 文件更大，保存和恢复更重 \\
分片检查点 & 按 rank 或 stage 分散保存的训练状态 & FSDP / ZeRO / 大规模训练 & 恢复与格式转换更复杂，需要统一工具链 \\
适配器权重 & LoRA/PEFT 增量参数 & 多任务适配器管理、灰度实验 & 只能依附特定底座与版本恢复含义 \\
\bottomrule
\end{tabularx}
\caption{检查点策略的核心不是“存没存”，而是“恢复时能回到哪个状态”}
\label{tab:checkpoint-granularity}
\end{table}

因此，平台期至少要把一句话写进训练规范里：\textbf{恢复训练所需的最小状态集是什么。} 一旦这句话说不清，实验中断后大家就会在“重新 warmup 还是继续衔接”“学习率到底该从哪接”“数据游标是不是又重跑了一遍”这些问题上不断返工。

\subsection{恢复演练要在故障发生前做，而不是在故障发生时学}
很多团队直到第一次训练中断，才真正去读自己的 checkpoint 逻辑。这往往已经太晚了。教材里必须把这条纪律写明白：\textbf{只会保存，不算具备恢复能力；只有定期演练过恢复，才算真的有恢复链。}

一次够用的恢复演练通常不复杂。做法可以很朴素：选一条中等成本实验，固定某个 step 人为打断；从最近 checkpoint 恢复；确认模型参数、优化器状态、学习率调度、数据进度和日志步数都能合理续上；再用一小段回放验证恢复前后的 loss 与吞吐没有明显异常漂移。只要这件事能稳定完成，团队才配说自己具备“自动恢复”。

这一步之所以重要，是因为分布式训练的恢复往往不是单点问题。可能是某个 rank 的分片没存全，可能是调度器状态没跟上，可能是数据游标重跑了，可能是 launcher 重启后 rank 到设备映射变了。平时不演练，这些问题会一起积到真正故障时爆出来。平台期真正值钱的，不是“我们理论上能恢复”，而是“我们知道从哪一步恢复、恢复后看哪几个信号验活”。

\subsection{成本核算和实验复现，本身就是大模型工程的一部分}
当团队进入多卡、多实验、多数据版本阶段，“会不会跑起来”已经不是唯一问题，“是否值得继续跑”同样关键。成本核算至少要覆盖下面四类指标：
\begin{enumerate}
  \item \textbf{GPU 小时}：一次实验用了多少卡时。
  \item \textbf{有效吞吐}：每秒真正处理了多少 token 或样本。
  \item \textbf{失败恢复成本}：中断后回到可继续状态要多久。
  \item \textbf{复现成本}：下次重跑时，是否能还原同样的数据、参数和代码状态。
\end{enumerate}

\begin{table}[H]
\centering
\small
\begin{tabularx}{\textwidth}{>{\raggedright\arraybackslash}p{2.9cm}>{\raggedright\arraybackslash}p{3.2cm}>{\raggedright\arraybackslash}X}
\toprule
\textbf{成本视角} & \textbf{应该记录什么} & \textbf{为什么重要} \\
\midrule
GPU 小时 & 世界规模、训练时长、设备型号 & 决定实验是否值得反复继续 \\
有效吞吐 & token/s、样本/s、利用率变化 & 判断扩卡是否真的带来收益 \\
恢复成本 & checkpoint 周期、中断恢复时间 & 决定平台期是否承受得起故障 \\
复现成本 & 数据版本、代码版本、随机种子、配置文件 & 决定结果能否被团队重新验证 \\
\bottomrule
\end{tabularx}
\caption{大规模工程里，算得清和跑得动同样重要}
\label{tab:training-cost-view}
\end{table}

\subsection{实验最小复现包：没有这包材料，结果就只是传说}
讲复现时，很多教材会停留在“记录随机种子”和“保存配置文件”。对平台期工程来说，这还不够。更可靠的做法，是给每次重要训练都生成一份\textbf{最小复现包}，让团队未来至少能回答：这次结果是用什么代码、什么数据、什么环境和什么启动方式跑出来的。

一份够用的最小复现包，通常至少包括：
\begin{itemize}
  \item 代码版本或提交号；
  \item 数据版本与切分快照；
  \item 训练配置、并行策略与启动命令；
  \item 环境摘要，例如驱动、CUDA、NCCL、核心依赖版本；
  \item checkpoint 指针、日志目录和关键评测结果。
\end{itemize}

这件事看上去像文书工作，实际上是在给未来的排障、比较和迁移留证据。没有这包材料，所谓“那次 8 卡效果特别好”很快就会沦为口口相传的印象，而不是团队能复用的工程资产。平台期真正拉开差距的，往往不是谁多跑了一次实验，而是谁能把一次实验变成后续十次决策的可信起点。

\subsection{把“卡忙不忙”变成可比数字：吞吐、Profiler 与有效算力}
旧稿里还花了不少篇幅讲 GPU 利用率、吞吐估计、PyTorch Profiler 和硬件诊断工具。新版如果只保留“要观测”这句原则，读者还是很难落到手上。因此，这里要补一个更可执行的判断：\textbf{训练观测至少要把 step time 拆开，并建立能跨实验比较的吞吐口径。}

对训练来说，最少要同时记录四类时间：
\begin{itemize}
  \item \textbf{纯计算时间}：前向、反向、优化器更新到底占了多少；
  \item \textbf{通信时间}：All-Reduce、All-Gather、Reduce-Scatter 等是否在拖后腿；
  \item \textbf{数据与 I/O 时间}：数据加载、预处理、checkpoint 落盘是否卡住流水；
  \item \textbf{空转时间}：等待其他 rank、等待节点、等待调度的时间到底有多少。
\end{itemize}

\begin{table}[H]
\centering
\small
\begin{tabularx}{\textwidth}{>{\raggedright\arraybackslash}p{2.8cm}>{\raggedright\arraybackslash}p{3.1cm}>{\raggedright\arraybackslash}X}
\toprule
\textbf{观测口径} & \textbf{最少记录什么} & \textbf{为什么它比单看 GPU 利用率更有用} \\
\midrule
集群吞吐 & token/s、样本/s、每卡吞吐 & 能直接比较扩卡前后到底有没有多干活 \\
Step 时间拆分 & 计算、通信、I/O、空转占比 & 能区分慢在模型、慢在网络还是慢在数据管线 \\
有效算力占比 & 结合理论 FLOPS 或 MFU/HFU 做粗对比 & 能判断当前更像算力没吃满，还是本质受带宽/通信约束 \\
Profiler / Trace & PyTorch Profiler、NCCL trace、系统级时间线 & 能把“感觉卡”变成具体哪一段卡的证据 \\
\bottomrule
\end{tabularx}
\caption{训练观测需要能跨实验比较，而不只是看一眼设备面板}
\label{tab:training-observability}
\end{table}

这里尤其要强调“吞吐估计法”的价值。很多团队会说“这次训练看起来更忙了”，但如果每秒处理的 token 没增加，或者扩到更多卡后每卡吞吐反而掉得厉害，那就说明新增资源并没有真正转化成有效训练。相反，只要统一记录每秒 token、每步时长和世界规模，就能很快看出扩卡收益到底有没有被通信和空转吃掉。

至于 FLOPS 比值、MFU 或 HFU 这类指标，不需要把它神秘化。它们更像一种{\heiti 横向比较工具}：不是要求你得到一个绝对精确的“理论最优值”，而是帮助你比较两套配置时，哪一套更接近把算力转成真实训练进展。对工程团队来说，这已经足够有价值。

\subsection{通信调优不是玄学：bucket、overlap 与 sampler 纪律都能吞掉收益}
当分布式训练扩到多卡、多机以后，很多“吞吐为什么没上去”的根因，已经不在模型结构，而在通信细节。旧稿里零散提到的 bucket size、重叠通信、样本切分纪律，这里需要收成一条明确的工程逻辑：\textbf{通信不是背景噪声，而是扩卡收益的主要税项。}

\begin{table}[H]
\centering
\small
\begin{tabularx}{\textwidth}{>{\raggedright\arraybackslash}p{2.8cm}>{\raggedright\arraybackslash}p{3cm}>{\raggedright\arraybackslash}X}
\toprule
\textbf{调节点} & \textbf{主要想改善什么} & \textbf{常见风险} \\
\midrule
梯度 bucket 大小 & 减少通信碎片或让 overlap 更容易生效 & bucket 太小会调用过多，太大又会推迟同步开始时机 \\
通信与计算 overlap & 让 All-Reduce / Reduce-Scatter 藏到反向过程中 & 若图结构或框架实现不合适，理论可重叠不代表真重叠 \\
长度分桶与 sampler 纪律 & 减少 rank 间 step 时长不一致 & 如果各卡样本长度差异过大，就会出现明显拖尾 rank \\
梯度累积频率 & 降低同步频率，换取更大有效 batch & 同步少了不等于更快，step time 和收敛口径都要重算 \\
\bottomrule
\end{tabularx}
\caption{通信调优的本质是减少无效等待，而不是追求参数越多越好}
\label{tab:communication-tuning}
\end{table}

这里尤其要强调 sampler 纪律。很多团队明明所有卡都在跑，却总有一张卡“最后才结束”。根因往往不是框架 bug，而是不同 rank 拿到的样本长度分布差得太远，导致某些卡总在处理更长的序列、更重的 packing 或更慢的预处理。只要这一点不收住，任何 bucket 或 overlap 调整都只能治标。

因此，平台期最好形成一个固定顺序：先看 rank 间 step time 是否均衡，再看通信是否和计算重叠，最后才去细调 bucket 大小和框架参数。这样做能避免把“数据不齐”误判成“通信太慢”，也能避免在错误层面上反复试配置。

\subsection{先定位哪一层卡住，再进入排障}
分布式训练里的排障不该从“猜框架有 bug”开始，而应从最便宜的检查做起：
\begin{enumerate}
  \item 先看是显存不够、吞吐低、启动失败还是恢复失败。
  \item 再看是计算瓶颈、通信瓶颈、I/O 瓶颈还是调度瓶颈。
  \item 最后才去怀疑更深的框架或环境问题。
\end{enumerate}

\begin{table}[H]
\centering
\small
\begin{tabularx}{\textwidth}{>{\raggedright\arraybackslash}p{3.2cm}>{\raggedright\arraybackslash}p{3.3cm}>{\raggedright\arraybackslash}X}
\toprule
\textbf{现象} & \textbf{最该先怀疑什么} & \textbf{第一检查动作} \\
\midrule
多卡扩展后吞吐几乎不涨 & 通信吞掉收益 & 看 All-Reduce 时间、桶大小、batch 是否太小 \\
首卡显存明显更高 & rank 0 额外逻辑过多 & 查日志、loss 汇总、保存与控制逻辑是否集中在一张卡 \\
训练经常中断后难恢复 & 检查点与恢复链不完整 & 核查是否同时保存了优化器、调度器和数据进度 \\
多机起不来或随机卡死 & NCCL / 网络 / 端口配置 & 先查连通性、环境变量、节点拓扑和启动参数 \\
GPU 利用率看着很高但训练很慢 & I/O 或流水等待被掩盖 & 结合 profiler、吞吐和数据加载时延一起看 \\
\bottomrule
\end{tabularx}
\caption{分布式训练排障先从最便宜的怀疑开始}
\label{tab:distributed-diagnosis}
\end{table}

\section{主案例推进}
回到企业知识助手。到本章时，它已经不再只是一个调用开源模型做问答的小系统，而是进入了需要持续改造的阶段：团队要维护多套适配器、要做继续预训练实验、要定期回灌新数据，还要评估更长上下文与更大底座模型。

\subsection{从原型期到平台期，升级路径不该跳步}
一个更现实的升级路径通常分三阶段：
\begin{enumerate}
  \item \textbf{原型期}：单机推理、少量 SFT / LoRA 实验，重点是把链路跑通。
  \item \textbf{成长期}：单机多卡训练、系统化评测、适配器版本治理，重点是把迭代跑稳。
  \item \textbf{平台期}：继续预训练、多任务适配器管理、统一日志与检查点、成本核算和自动恢复，重点是把组织能力搭起来。
\end{enumerate}

很多项目真正的问题不是“没有上 3D 并行”，而是原型期就试图使用平台期复杂度，导致链路在尚未稳定时过早爆炸。

\subsection{企业助手在平台期的第一批分布式动作}
对我们的主案例而言，更合理的顺序通常是：
\begin{enumerate}
  \item 先用单机或单节点把 LoRA / QLoRA 和评测闭环跑稳。
  \item 单卡装不下时，先尝试 FSDP 或 ZeRO 一类分片方案。
  \item 只有在继续预训练或更大模型明确需要时，再评估 TP 或 PP。
  \item 等多类任务同时进入训练管线后，再建设统一实验平台、适配器仓库和自动恢复机制。
\end{enumerate}

这条顺序背后有一个非常朴素的原则：\textbf{复杂度应该随着需求增长，而不是随着技术好奇心增长。}

\subsection{平台期的第一次恢复演练，比第一次扩到更多卡更值钱}
回到企业知识助手这个主案例，团队进入平台期后，最诱人的动作往往是“扩更多卡、训更大模型、同时开更多实验”。但从长期收益看，更值得优先做的一步，通常是\textbf{把恢复演练、配置归档和实验账本先跑通。}

原因很简单。企业助手的继续预训练和多任务适配实验一旦开始变贵，真正拖垮团队节奏的，往往不是单次吞吐不够高，而是一次中断后恢复不了、一次版本变更后说不清差异、一次框架迁移后旧实验全断档。也就是说，平台期的第一张门票不是“会开更多卡”，而是“训练出了事还能把局接住”。

因此，对主案例更现实的一组优先级是：先统一启动模板，再统一 checkpoint 规范，再做一次恢复验活，最后才逐步扩 world size。只要这条顺序立住，后面不论是继续预训练、适配器多任务训练，还是更重的分布式方案，团队都不是在裸奔。

\subsection{平台期必须补上的治理清单}
一旦企业知识助手进入平台期，下面这些能力会比“会不会写 distributed script”更值钱：
\begin{itemize}
  \item \textbf{数据版本治理}：知道每个实验吃的是什么样本集合。
  \item \textbf{配置可追踪}：知道每次训练用了什么并行策略、学习率、batch 和恢复点。
  \item \textbf{适配器资产管理}：知道哪套适配器对应哪个任务、哪批数据和哪次评测。
  \item \textbf{自动恢复链}：知道节点失败后谁来重启、从哪里恢复、恢复后怎样验活。
  \item \textbf{成本账本}：知道哪类实验长期看收益最高，哪类实验只是烧算力。
\end{itemize}

\subsection{把前十二章重新串起来：先形成第一条主线，再进入系统深化}
走到这里，前十二章已经可以先被看成一条闭合的基础主线，而不是十二个彼此独立的专题：
\begin{enumerate}
  \item 第 1--2 章定义模型是什么、底层信息流怎么工作；
  \item 第 3--6 章解决“如何把模型变成可控的任务执行者”；
  \item 第 7--8 章解决“如何把企业知识稳定接进系统”；
  \item 第 9--11 章解决“如何知道系统变好了，以及如何把它稳定地发出去”；
  \item 第 12 章再把这些能力抬到组织级工程，回答“如何持续、大规模、可恢复地迭代”。
\end{enumerate}

这个回看非常重要，因为平台期最容易发生的失败，不是单点技术不会，而是团队忘了这些能力本来就应该首尾相接。没有模型与任务边界，后面的训练会盲；没有知识治理，后面的评测会飘；没有评测和门禁，后面的分布式训练只会更快地产生更难追责的版本。

但这还不是全书终点。更准确地说，前十二章到这里完成的是“模型到底座工程”的第一条主线：从原理理解、任务控制、数据与微调、RAG、评测、对齐、推理优化，一路走到训练平台和恢复治理。只有这条主线站稳，后面的四章才有意义，因为系统接下来会进入另一层复杂度：它不只要会答、会训、会评，还要会跨工具执行、会处理更深推理、会守住高权限环境下的安全边界，并且能在生产环境里长期运行。

\subsection{平台期真正交付的，不是更大模型，而是更可靠的迭代能力}
很多团队把平台期想成“终于能训更大的模型了”。这当然是平台期的一部分，但不是最重要的那部分。对企业知识助手来说，平台期真正交付的成果更像下面这张表：

\begin{table}[H]
\centering
\small
\begin{tabularx}{\textwidth}{>{\raggedright\arraybackslash}p{2.8cm}>{\raggedright\arraybackslash}p{3cm}>{\raggedright\arraybackslash}X}
\toprule
\textbf{平台期能力} & \textbf{表面上看像什么} & \textbf{它真正保障了什么} \\
\midrule
统一训练与恢复模板 & 一套脚本、一套配置规范 & 故障发生时，团队能接住训练，不至于从头来过 \\
数据与适配器版本治理 & 一堆实验记录与资产编号 & 每次收益和退化都能找到来源，不再靠记忆对账 \\
评测与门禁接入训练链 & 训练完自动跑回归和关键切片 & 版本不只会被训练出来，还会被认真拦住或放行 \\
成本账本与实验优先级 & GPU 小时、吞吐、失败率报表 & 团队知道哪些方向值得持续投，哪些只是昂贵幻觉 \\
\bottomrule
\end{tabularx}
\caption{平台期最重要的交付物，不只是更大的模型，而是更可靠的迭代机制}
\label{tab:platform-delivery}
\end{table}

这张表背后的闭环感很强：前面十一章教会读者“问题该怎么分层、知识该怎么接、模型该怎么调、效果该怎么评”；第十二章则把这些局部能力收束成一套组织可以长期复用的工程节奏。只有到了这里，“企业知识助手”才真正从一个项目变成一种能力。

\subsection{平台期验收清单：什么时候说明训练与治理底座已经站稳}
如果要判断平台期有没有真正站稳，可以用下面这张自查表来看训练与治理底座是否已经形成：
\begin{itemize}
  \item \textbf{模型边界清楚}：团队知道哪些问题该修模型、该修任务、该修系统。
  \item \textbf{知识链路清楚}：团队知道答案来自哪类文档、哪类工具、哪条权限边界。
  \item \textbf{评测门禁清楚}：团队知道哪些切片能拦版本、哪些只是观察项。
  \item \textbf{训练治理清楚}：团队知道实验吃了哪批数据、用了哪套配置、从哪里恢复。
  \item \textbf{成本责任清楚}：团队知道哪类优化真的带来了长期收益，哪类只是算力消耗。
\end{itemize}

只要这五件事能被明确回答，就说明企业知识助手已经具备了继续往后走的底座。因为后面的 Agent、多步推理、安全治理和生产运维，都会默认训练、知识、评测和恢复链已经不是黑盒；如果这层底座没站稳，后面四章只会把复杂度叠得更高，而不会把系统真正做强。

\section{常见坑与工程取舍}
\begin{itemize}
  \item \textbf{单卡还能跑时就急着上复杂并行}：引入额外复杂度，却未必有实际收益。
  \item \textbf{忽略通信成本}：多卡不一定更快，网络条件差时收益会被快速吃掉。
  \item \textbf{把 TP / PP 当成“更高级”的默认方案}：很多时候更基础的分片方案已经够用。
  \item \textbf{没有检查点策略就长时间训练}：只要中断一次，代价就会非常高。
  \item \textbf{只看 GPU 利用率，不看实际吞吐}：看起来“卡很忙”不等于训练真的高效。
  \item \textbf{实验记录不完整}：下周无法复现本周结果，平台期就会失去工程意义。
\end{itemize}

再补一份更贴近一线的排障最小手册：
\begin{itemize}
  \item \textbf{多机启动失败}：先查环境变量、端口、主节点配置、节点连通性与 NCCL 设定。
  \item \textbf{吞吐提不上去}：先查 batch 是否过小、数据加载是否堵塞、通信是否压过计算。
  \item \textbf{首卡压力异常大}：先查 loss 聚合、日志打印、保存逻辑和额外控制流是否都落在 rank 0。
  \item \textbf{checkpoint 特别慢}：先查存储带宽、保存频率与保存内容，而不是先怪框架。
  \item \textbf{恢复训练后 loss 异常}：先查优化器状态、学习率调度状态和数据游标是否一起恢复。
\end{itemize}

\section{本章练习}
\begin{exercise}[并行策略选择]
如果模型刚好单卡放不下，但团队的多卡通信条件一般，且当前主要目标是先把训练跑稳，此时通常应优先考虑哪类方案？为什么不建议一上来就用张量并行或流水线并行？
\end{exercise}

\begin{solution}
通常应优先考虑{\heiti 数据并行配合参数分片}，例如 FSDP、ZeRO 这类方案。因为当前的主矛盾首先是{\heiti 常驻显存不够}，而不是单层算子必须切开，也不是模型深度必须做阶段调度。张量并行和流水线并行当然能解决更重的问题，但它们会引入更高的通信和调度复杂度。若网络条件一般、团队还在追求先跑稳，那么先用分片方案往往是更好的第一步。
\end{solution}

\begin{exercise}[训练故障归因]
为什么大规模训练里的“日志、检查点、自动恢复和成本统计”不是附属工作，而是核心基础设施？请从故障恢复和实验决策两个角度回答。
\end{exercise}

\begin{solution}
从{\heiti 故障恢复}角度看，训练时间一长，中断几乎不可避免。没有日志就难以知道故障发生在哪一层，没有检查点就难以回到可继续状态，没有自动恢复就需要人长期盯守，整体训练效率会被中断成本严重拖垮。从{\heiti 实验决策}角度看，没有成本统计与版本记录，就无法判断某次实验到底值不值得继续，也无法确认“这次收益”究竟来自数据、超参数、并行策略还是偶然因素。平台期的大模型工程，本质上就是把训练过程变成{\heiti 可追踪、可恢复、可核算}的系统工程。
\end{solution}

\section{延伸阅读}
\begin{itemize}
  \item \textbf{PyTorch Distributed 与 FSDP 官方文档}：适合理解原生多卡训练和参数分片的基本接口与约束。
  \item \textbf{DeepSpeed 官方文档}：适合理解 ZeRO、offload、训练配置治理等大规模工程能力。
  \item \textbf{Megatron-LM 及相关技术报告}：适合理解 TP、PP 与 3D 并行在超大模型里的组合方式。
  \item \textbf{NCCL、集群拓扑与 profiler 实践资料}：适合理解为什么“能跑”“跑得快”“跑得稳”是三件不同的事。
\end{itemize}

\section{小结}
本章把分布式训练从“技术名词表”收束成了一套工程判断：先看瓶颈属于显存、层内计算、模型深度还是通信，再决定用 DDP、FSDP、ZeRO、TP、PP 还是更复杂的组合；先把检查点、恢复、日志与成本账本建起来，再谈大规模稳定迭代。对企业知识助手来说，平台期真正决定成败的，往往不是某一种并行技巧，而是整套训练治理体系是否站得住。

\section{下一章过渡}
当训练平台和恢复链站稳以后，系统复杂度的下一次跃迁，往往不再来自“再多几张卡”，而来自“再多几个动作”。企业知识助手不只要会答、会训，还要开始跨检索、日历、审批和邮件系统执行真实流程。下一章就转向 Agent 与工具编排：在已有最小 tool calling 基础上，把模型升级为可规划、可审批、可回滚的多步执行者。
